% GENERATED FILE. Edit canonical lessons, then run npm run generate:books.
% canonical-source-hash: fnv1a64:bce48bba0f215694
% canonical-lessons: ZH-C298-huozai, ZH-C298-jiuhuo, ZH-W298-jiu-save, ZH-C298-xiaofangyuan, ZH-C298-han, ZH-C298-jiuming

\chapter{A Fire}
\label{ch:zh-a-fire}
\hlchaptermodality{298}

\begin{chapteropening}
\textbf{By the end of this chapter:} \emph{I can talk about a fire, the firefighters and someone shouting for help.}

Complete the last lesson of chapter 298: I can talk about a fire, the firefighters and someone shouting for help.
\end{chapteropening}

\section[huǒzāi]{huǒzāi — big fire; a fire that destroys buildings}

\label{lesson:ZH-C298-huozai}



\begin{quote}
Before the new one: say the Mandarin for what comes of it, then the Mandarin for to deal with a problem.
\end{quote}

\subsection*{You'll want to know: huǒzāi}
\emph{huǒzāi} — \textquotedblleft{}big fire; a fire that destroys buildings\textquotedblright{}.

\begin{quote}
\emph{huǒzāi de yuányīn shì shénme} — What caused the fire?
\end{quote}

Its characters come later; for now it is a word for the ear.

\begin{cousinweb}
\emph{huǒzāi} is built from \zh{火} 'fire' and \emph{zāi} 'a disaster', both inherited from Middle Chinese, \emph{xwaX} and \emph{tsoj} (Guangyun, in Baxter's transcription), which lists the second as a spelling of the word for fire sent from heaven.
\end{cousinweb}

\begin{grammarlens}[title={What you've built}]
One more word for telling what happened.
\end{grammarlens}

\subsection*{Your turn}
\begin{itemize}
\raggedright
  \item \emph{Say it:} \emph{huǒzāi}
  \item \emph{Say it:} \emph{huǒzāi}, once more
  \item \emph{Say it:} \emph{huǒzāi de yuányīn shì shénme} — What caused the fire?
  \item \emph{Recall:} say the Mandarin for what comes of it, then the Mandarin for to deal with a problem, then say \emph{huǒzāi} again
\end{itemize}

\begin{tcolorbox}[breakable,colback=teal!4,colframe=teal!35!black,title={Before you move on}]
What does huǒzāi mean? (\textquotedblleft{}Big fire; a fire that destroys buildings\textquotedblright{}.) Say it once more.
\end{tcolorbox}
\section[\texorpdfstring{\zh{救火} (jiùhuǒ)}{救火 (jiùhuǒ)}]{\zh{救火} (jiùhuǒ) — to fight a fire; to put out a fire}

\label{lesson:ZH-C298-jiuhuo}



\begin{quote}
Before the new one: say the Mandarin for to deal with a problem, then the Mandarin for a fire that destroys buildings.
\end{quote}

\subsection*{You'll want to know: \zh{救火}}
\textbf{\zh{救火}} — \emph{jiùhuǒ} — \textquotedblleft{}to fight a fire; to put out a fire\textquotedblright{}.

\begin{quote}
\textbf{\zh{大家都来救火了}} \emph{dàjiā dōu lái jiùhuǒ le} — Everyone came to fight the fire.
\end{quote}

The new character, in pieces: \zh{救} — eleven strokes in two pieces: on the left seven strokes, a bar, a vertical with a hook, a dot and a rising stroke on its left, a short sweep and a dot on its right, and a dot at the top right; on the right the striking piece of \zh{故}, four strokes. The next lesson puts it in your hand.

\begin{cousinweb}
\textbf{\zh{救火}} (\emph{jiùhuǒ}) is built from \zh{救} 'to save' and \zh{火} 'fire', both inherited from Middle Chinese, \emph{kjuwH} and \emph{xwaX} (Guangyun, in Baxter's transcription), which glosses the first 'to protect' and 'to stop': to stop a fire.
\end{cousinweb}

\begin{grammarlens}[title={What you've built}]
One more word for telling what happened.
\end{grammarlens}

\subsection*{Your turn}
\begin{itemize}
\raggedright
  \item \emph{Say it:} \emph{jiùhuǒ}
  \item \emph{Say it:} \emph{jiùhuǒ}, once more
  \item \emph{Say it:} \emph{dàjiā dōu lái jiùhuǒ le} — Everyone came to fight the fire.
  \item \emph{Recall:} say the Mandarin for to deal with a problem, then the Mandarin for a fire that destroys buildings, then say \emph{jiùhuǒ} again
\end{itemize}

\begin{tcolorbox}[breakable,colback=teal!4,colframe=teal!35!black,title={Before you move on}]
What does \zh{救火} mean? (\textquotedblleft{}To fight a fire; to put out a fire\textquotedblright{}.) Say it once more.
\end{tcolorbox}
\section[\texorpdfstring{\zh{救} (jiù)}{救 (jiù)}]{\zh{救} — a hooked post with small strokes, then the striking piece}

\label{lesson:ZH-W298-jiu-save}



\begin{quote}
The word first, then the character inside it.

\begin{itemize}
\raggedright
  \item \emph{Recall:} say \emph{jiùhuǒ} — \textbf{R1}, one lesson back, before it decays
\end{itemize}

\begin{quote}
\zh{故} \zh{火}
\end{quote}

\emph{Read it:} each one aloud, with nothing said first — \emph{gù}, \emph{huǒ}
\end{quote}

\subsection*{Script — a hooked post with small strokes, then the striking piece}
\hlblockfigure{\detokenize{figures/ZH-W298-jiu-save-filmstrip.pdf}}{How \zh{救} is written, stroke by stroke}

\begin{quote}
\zh{救}
\end{quote}

\textbf{Two pieces, left before right.} Count the pieces, not the strokes.

\textbf{The left piece, seven strokes.} A bar. A vertical down through it that hooks to the left at the foot. A dot on its left, then a stroke rising to the right below the dot. A short sweep down to the left on its right, then a short stroke falling to the right under that. Last, a dot high on the right, above the end of the bar.

\textbf{The right piece, four strokes.} The striking piece of \zh{故}: a sweep to the left, a bar, a long sweep to the lower left and a long stroke falling to the right.

7 + 4 = eleven.

\subsection*{Your turn}
\begin{itemize}
\raggedright
  \item \emph{Write it:} \zh{救}, all eleven, with the model in front of you
  \item \emph{Cover:} write it again from memory, and count the pieces, not the strokes
\end{itemize}

\begin{tcolorbox}[breakable,colback=teal!4,colframe=teal!35!black,title={Before you move on}]
How many strokes in \zh{救}? (\textbf{Eleven}.) Which stroke of the left piece comes last? (\textbf{The dot high on the right}.) Where else have you made the right piece? (\textbf{In \zh{故} and \zh{做}}.)

Source: \href{https://raw.githubusercontent.com/chanind/hanzi-writer-data/68d10a4b21150cae5e1ebbd223eed289cf32d90c/data/\%E6\%95\%91.json}{Hanzi Writer Data \zh{救}, pinned stroke paths and medians}.
\end{tcolorbox}
\section[xiāofángyuán]{xiāofángyuán — firefighter}

\label{lesson:ZH-C298-xiaofangyuan}



\begin{quote}
Before the new one: say the Mandarin for a fire that destroys buildings, then the Mandarin for to put out a fire.
\end{quote}

\subsection*{You'll want to know: xiāofángyuán}
\emph{xiāofángyuán} — \textquotedblleft{}firefighter\textquotedblright{}.

\begin{quote}
\emph{xiāofángyuán hěn kuài jiù lái le} — The firefighters were there in no time.
\end{quote}

Its characters come later; for now it is a word for the ear.

\begin{cousinweb}
\emph{xiāofángyuán} is built from \emph{xiāo} 'to put out', \emph{fáng} 'to guard against' and \emph{yuán} 'a member', all inherited from Middle Chinese, \emph{sjew}, \emph{bjang} and \emph{hjwen} (Guangyun, in Baxter's transcription), which glosses the first 'to put out', the second 'to keep off' and the third 'a count of things': one of those who put out fires.
\end{cousinweb}

\begin{grammarlens}[title={What you've built}]
One more word for telling what happened.
\end{grammarlens}

\subsection*{Your turn}
\begin{itemize}
\raggedright
  \item \emph{Say it:} \emph{xiāofángyuán}
  \item \emph{Say it:} \emph{xiāofángyuán}, once more
  \item \emph{Say it:} \emph{xiāofángyuán hěn kuài jiù lái le} — The firefighters were there in no time.
  \item \emph{Recall:} say the Mandarin for a fire that destroys buildings, then the Mandarin for to put out a fire, then say \emph{xiāofángyuán} again
\end{itemize}

\begin{tcolorbox}[breakable,colback=teal!4,colframe=teal!35!black,title={Before you move on}]
What does xiāofángyuán mean? (\textquotedblleft{}Firefighter\textquotedblright{}.) Say it once more.
\end{tcolorbox}
\section[hǎn]{hǎn — to shout; to call out}

\label{lesson:ZH-C298-han}



\begin{quote}
Before the new one: say the Mandarin for to put out a fire, then the Mandarin for a firefighter.
\end{quote}

\subsection*{You'll want to know: hǎn}
\emph{hǎn} — \textquotedblleft{}to shout; to call out\textquotedblright{}.

\begin{quote}
\emph{wǒ tīngdào yǒu rén zài hǎn} — I heard someone shouting.
\end{quote}

Its characters come later; for now it is a word for the ear.

\begin{cousinweb}
\emph{hǎn} is inherited from Middle Chinese \emph{xamX} (Guangyun, in Baxter's transcription), which glosses it 'a sound'; the Guangyun also lists it under two other readings, \emph{heamX} and \emph{xeamX}, both glossed as a sound or a cry.
\end{cousinweb}

\begin{grammarlens}[title={What you've built}]
One more word for telling what happened.
\end{grammarlens}

\subsection*{Your turn}
\begin{itemize}
\raggedright
  \item \emph{Say it:} \emph{hǎn}
  \item \emph{Say it:} \emph{hǎn}, once more
  \item \emph{Say it:} \emph{wǒ tīngdào yǒu rén zài hǎn} — I heard someone shouting.
  \item \emph{Recall:} say the Mandarin for to put out a fire, then the Mandarin for a firefighter, then say \emph{hǎn} again
\end{itemize}

\begin{tcolorbox}[breakable,colback=teal!4,colframe=teal!35!black,title={Before you move on}]
What does hǎn mean? (\textquotedblleft{}To shout; to call out\textquotedblright{}.) Say it once more.
\end{tcolorbox}
\section[jiùmìng]{jiùmìng — help!; the cry for help}

\label{lesson:ZH-C298-jiuming}



\begin{quote}
Before the new one: say the Mandarin for a firefighter, then the Mandarin for to shout loudly.
\end{quote}

\subsection*{You'll want to know: jiùmìng}
\emph{jiùmìng} — \textquotedblleft{}help!; the cry for help\textquotedblright{}.

\begin{quote}
\emph{tā dà hǎn jiùmìng} — He shouted for help.
\end{quote}

Its characters come later; for now it is a word for the ear.

\begin{cousinweb}
\emph{jiùmìng} is built from \zh{救} 'to save' and \emph{mìng}, both inherited from Middle Chinese, \emph{kjuwH} and \emph{mjaengH} (Guangyun, in Baxter's transcription), which glosses the second 'to order' and 'to send'; its sense 'a life' is not given there.
\end{cousinweb}

\begin{grammarlens}[title={What you've built}]
One more word for telling what happened.
\end{grammarlens}

\subsection*{Your turn}
\begin{itemize}
\raggedright
  \item \emph{Say it:} \emph{jiùmìng}
  \item \emph{Say it:} \emph{jiùmìng}, once more
  \item \emph{Say it:} \emph{tā dà hǎn jiùmìng} — He shouted for help.
  \item \emph{Recall:} say the Mandarin for a firefighter, then the Mandarin for to shout loudly, then say \emph{jiùmìng} again
\end{itemize}

\begin{tcolorbox}[breakable,colback=teal!4,colframe=teal!35!black,title={Before you move on}]
What does jiùmìng mean? (\textquotedblleft{}Help!; the cry for help\textquotedblright{}.) Say it once more.
\end{tcolorbox}
